\chapter{Integration by Substitution}

In the previous chapter, the Fundamental Theorem of Calculus showed that
computing an accumulated area $\int_a^b h(x)\,\dd x$ does
not require adding up rectangles. It only requires finding
an antiderivative---a function $F$ whose derivative is $h$---and evaluating
$F(b) - F(a)$.

The catch is finding that $F$.

We can readily differentiate any standard elementary function: combinations of
powers, exponentials, logarithms, and trigonometric functions assembled through sums,
products, quotients, and compositions. The chain rule, in particular, makes compositions
routine: the derivative of $f(g(x))$ is $f'(g(x))\,g'(x)$, where the factor $g'(x)$ accounts
for the rate of change of the inner function.

Integration is fundamentally harder, and there is a real reason for it. Differentiation is
a \emph{procedure}: follow the rules and the answer falls out mechanically. Integration is a
\emph{search}. To integrate an \emph{integrand} $h(x)$ -- the function sitting inside the
integral waiting to be integrated -- is to find some $F$ whose derivative is $h$, and
there is no mechanical rule that always produces it. So we develop a toolkit of
techniques, each one a method for \emph{recognizing} an integrand as
the derivative of a known function.

This chapter is about the first and most important tool in that box:
\textbf{substitution}, also called $u$-substitution. Here is the core idea of
the chapter in one sentence:

\begin{focusbox}
Substitution is the chain rule read backwards. The chain rule tells you how the
derivative of a composition picks up an extra factor from the inside; substitution
teaches your eye to spot that extra factor in an integrand and to undo the
composition.
\end{focusbox}

We will build the idea from a picture, prove that it works (it is a one-line
consequence of the chain rule and the Fundamental Theorem), learn how to choose
the substitution, sort out the delicate business of what happens to $\dd x$, and
handle the definite-integral case where the limits themselves get carried along.
Then we will work a long sequence of examples, from the obvious to the genuinely
sneaky, and pin down what the symbol-shuffling really means.

\chapterline

\section{The pattern that starts it all}
\label{ch22:the-pattern}

Suppose we want to integrate
\[
  \int 2x\cos\!\left(x^2\right)\,\dd x.
\]
Stare at it. There is a composition in there, $\cos(x^2)$, whose inside is $x^2$.
And sitting right next to it is $2x$ -- which is the exact derivative of that
inside. That matching factor is our signal.

Because you know the chain rule, you can run a quick experiment. What is the
derivative of $\sin(x^2)$? The outside function is $\sin$, the inside is $x^2$, so
\[
  \ddx{\sin(x^2)} = \cos(x^2)\cdot \underbrace{2x}_{\text{from the inside}}
                  = 2x\cos(x^2).
\]
That is exactly our integrand. So
\[
  \int 2x\cos\!\left(x^2\right)\,\dd x = \sin(x^2) + C,
\]
and you can check it by differentiating the answer. The whole trick worked
because the integrand had the shape
\[
  \big(\text{outside}'\ \text{of the inside}\big)\times\big(\text{derivative of the inside}\big),
\]
which is precisely the shape the chain rule manufactures. Substitution is the
discipline of \emph{seeing} this shape even when it is disguised, and of having a
reliable bookkeeping system to exploit it.

\begin{intuitionbox}
Think of the chain rule as a machine that, when it differentiates $F(g(x))$, spits
out two pieces multiplied together: the ``outer derivative'' $F'(g(x))$ and the
``inner derivative'' $g'(x)$. If you are ever handed an integrand that looks like
those two pieces multiplied together, you are looking at the output of that machine
-- and you can run the machine in reverse to recover $F(g(x))$.
\end{intuitionbox}

\begin{whybox}
The mechanism, stated cleanly. The chain rule says
\[
  \ddx{F(g(x))} = F'(g(x))\,g'(x).
\]
Read that equation from right to left. It says: the antiderivative of
$F'(g(x))\,g'(x)$ is $F(g(x))$. In integral notation,
\[
  \int F'(g(x))\,g'(x)\,\dd x = F(g(x)) + C.
\]
Substitution is nothing more than a clean way of \emph{recognizing} an integrand as
having the form $F'(g(x))\,g'(x)$ and then writing down $F(g(x))$. Every $u$-sub you
will ever do is an instance of this single equation. The only skill is learning to
see the pattern and to manage the bookkeeping when it is partly hidden.
\end{whybox}

\subsection{Naming the inside}

The mental move that turns this from a lucky guess into a method is to give the
inside function a name. Call the inside $u$. In our example $u = x^2$. Then its
derivative is $\dv{u}{x} = 2x$, which physicists and Leibniz alike happily rearrange
into the suggestive form
\[
  \dd u = 2x\,\dd x.
\]
Look what that does to the integral. Wherever we see $x^2$ we write $u$; wherever we
see $2x\,\dd x$ we write $\dd u$:
\[
  \int \cos\!\big(\underbrace{x^2}_{u}\big)\,\underbrace{2x\,\dd x}_{\dd u}
   = \int \cos u\,\dd u = \sin u + C = \sin(x^2)+C.
\]
The hard integral in $x$ became a trivial integral in $u$. \emph{That} is the
payoff: a good substitution replaces a tangled integrand with a simple one, does
the easy integral, and then translates back. We will spend the rest of the chapter
making this precise and learning to do it when the pattern is buried.

\begin{recoverybox}
\textbf{Pause and think.} Look at $\int 3x^2 e^{x^3}\,\dd x$. Before writing anything down, spot the inside function and its derivative. What is $u$? What is $\dd u$? Can you see the antiderivative in your head in five seconds?

\emph{Answer:} Let $u = x^3$, so $\dd u = 3x^2\,\dd x$. The integrand is simply $e^u\,\dd u$, which integrates immediately to $e^u + C = e^{x^3} + C$.
\end{recoverybox}

\begin{figure}[H]
\centering
\begin{tikzpicture}
\begin{axis}[
  width=12cm, height=6.4cm,
  xlabel={$x$}, ylabel={$y$},
  domain=-2.4:2.4, samples=160,
  axis lines=middle,
  xtick={-1.5,-1,1,1.5},
  enlargelimits=true,
  legend style={at={(0.5,1.03)}, anchor=south, legend columns=2, font=\small, /tikz/every even column/.append style={column sep=14pt}},
  legend cell align=left,
]
\addplot[curveblue, thick] {sin(deg(x^2))};
\addlegendentry{$F(x)=\sin(x^2)$}
\addplot[tangentred, thick, dashed] {2*x*cos(deg(x^2))};
\addlegendentry{$F'(x)=2x\cos(x^2)$}
\end{axis}
\end{tikzpicture}
\caption{The antiderivative $F(x)=\sin(x^2)$ (solid) and its derivative
$F'(x)=2x\cos(x^2)$ (dashed). The dashed curve is the integrand we were handed; the
solid curve is what substitution recovers. Notice the dashed curve reaches its peaks and troughs exactly where the solid curve is rising or falling fastest -- that is the
chain rule made visible.}
\label{ch22:fig-pattern}
\end{figure}

\section{The theorem, with proof}
\label{ch22:theorem}

Intuition is wonderful, but we promised rigor. Here is substitution stated as a
theorem and proved, both for indefinite and definite integrals. The proof is
short because all the real work was done when we proved the chain rule and the
Fundamental Theorem of Calculus; substitution just bolts them together.

\begin{theorem}[Substitution, indefinite form]
\label{ch22:thm-indef}
Let $g$ be differentiable on an interval and let $F$ be an antiderivative of a
continuous function $f$ on the range of $g$, so that $F' = f$. Then
\[
  \int f\!\big(g(x)\big)\,g'(x)\,\dd x = F\!\big(g(x)\big) + C.
\]
\end{theorem}

\begin{proof}
We do not need to integrate anything. We only need to check that the proposed
answer, $F(g(x))$, has the right derivative -- because an antiderivative is, by
definition, any function whose derivative is the integrand. Apply the chain rule
to $F(g(x))$:
\[
  \ddx{F\!\big(g(x)\big)} = F'\!\big(g(x)\big)\,g'(x) = f\!\big(g(x)\big)\,g'(x).
\]
The first equality is the chain rule; the second uses $F'=f$. So $F(g(x))$ is
indeed an antiderivative of $f(g(x))\,g'(x)$, and every antiderivative differs from
it by a constant. That is exactly the claim.
\end{proof}

That is the whole content. The Leibniz notation $\dd u = g'(x)\,\dd x$ is a
mnemonic device for this theorem -- a way of doing the substitution on paper so
that the pattern is impossible to miss. But the theorem stands on its own without
any talk of ``infinitesimals.''

\begin{theorem}[Substitution, definite form]
\label{ch22:thm-def}
With the same hypotheses, and assuming $g'$ is continuous on $[a,b]$,
\[
  \int_a^b f\!\big(g(x)\big)\,g'(x)\,\dd x
   = \int_{g(a)}^{g(b)} f(u)\,\dd u.
\]
\end{theorem}

\begin{proof}
The Fundamental Theorem of Calculus evaluates the left-hand side using the
antiderivative we just found:
\[
  \int_a^b f\!\big(g(x)\big)\,g'(x)\,\dd x
   = \eval{F\!\big(g(x)\big)}{x=a}^{x=b}
   = F\!\big(g(b)\big) - F\!\big(g(a)\big).
\]
Now look at the right-hand side. Since $F'=f$, the Fundamental Theorem also gives
\[
  \int_{g(a)}^{g(b)} f(u)\,\dd u = \eval{F(u)}{u=g(a)}^{u=g(b)}
   = F\!\big(g(b)\big) - F\!\big(g(a)\big).
\]
Both sides equal the same difference $F(g(b)) - F(g(a))$, so they are equal.
\end{proof}

\begin{focusbox}
The definite-form theorem contains an essential instruction hiding in the limits of
integration: when you change the variable from $x$ to $u = g(x)$, you must also
change the limits. The new lower limit is $u = g(a)$ and the new upper limit is
$u = g(b)$. The endpoints travel with the substitution. Forgetting this is the most
common mistake in the whole chapter, and we will hammer it later.
\end{focusbox}

\begin{historybox}
Newton and Leibniz developed early calculus to solve \emph{quadrature}---finding
areas bounded by curves. Newton worked with \emph{fluxions}, viewing changing
variables as shifts in time or scale. Leibniz introduced the notation still in
use: the elongated sum $\int$ and the differential element $\dd x$.

Leibniz's notation makes $\dd u = g'(x)\,\dd x$ perform a coordinate change
directly, scaling differential widths like ordinary fractions. Euler developed
this systematic bookkeeping in his 1768 treatise \emph{Institutiones Calculi
Integralis}, interpreting substitution geometrically: $g'(x)$ acts as a local
stretch factor scaling the $x$-axis onto the $u$-axis so that area is preserved.

For over a century, the foundations remained informal. Early mathematicians
manipulated $\dd x$ and $\dd u$ as infinitesimals. In 1734, Bishop George Berkeley
objected in \emph{The Analyst} to dividing by small increments as nonzero quantities
and then setting those same increments to zero to finish the calculation.
Treating $\dd u / \dd x$ as a fraction gave correct results, but lacked a rigorous
foundation until Cauchy (1823) and Weierstrass formalized calculus using limits,
confirming that substitution is the Chain Rule integrated over an interval.
\end{historybox}

\section{What $\dd x$ is really doing}
\label{ch22:what-dx-does}

You have been told to ``cancel the $\dd x$,'' and it works, and that should bother
you a little. What is actually going on when we write $\dd u = 2x\,\dd x$ and then
treat the $\dd$'s like algebra?

\begin{happeningbox}
The $\dd x$ inside an integral is not a number being
multiplied -- it is a \emph{bookkeeping tag} that records which variable you are
integrating against, and, more deeply, it records how a tiny step in your variable
gets stretched into a tiny step in area.

When you substitute $u = g(x)$, a small change $\Delta x$ in $x$ produces a small
change $\Delta u \approx g'(x)\,\Delta x$ in $u$. The factor $g'(x)$ is the local
\emph{stretch factor}: it says how much the $u$-axis is being compressed or
expanded relative to the $x$-axis near that point. The integral
$\int f(g(x))\,g'(x)\,\dd x$ is adding up little rectangles of height $f(g(x))$ and
width $\dd x$. To rewrite this as an integral in $u$, each width $\dd x$ must be
re-measured in $u$-units, and the conversion factor is exactly $\dd u = g'(x)\,\dd
x$. So the ``cancellation'' is not a trick: $g'(x)\,\dd x$ genuinely \emph{is} the
width of the corresponding sliver measured along the $u$-axis. The equation
$\dd u = g'(x)\,\dd x$ is a statement about how widths transform, and Theorem
\ref{ch22:thm-def} is the rigorous version of that statement.
\end{happeningbox}

\begin{figure}[H]
\centering
\begin{tikzpicture}
% Top lane: x-ruler (uniform spacing dx)
\draw[thick, ->] (0.5,3.2) -- (11.5,3.2) node[right, font=\small] {$x$};
\node[font=\footnotesize\bfseries, sourceblue, anchor=south west] at (0.5,4.1) {Top Lane: Uniform $x$-ruler (constant width $\dd x$)};

% Grid ticks on x
\foreach \x/\label in {1.5/{$x_0{=}1$}, 3.5/{$x_1{=}2$}, 5.5/{$x_2{=}3$}, 7.5/{$x_3{=}4$}, 9.5/{$x_4{=}5$}} {
  \draw[thick, sourceblue] (\x,3.05) -- (\x,3.35) node[above, font=\footnotesize] {\label};
}
% Slices on x
\draw[|-|, thick, sourceblue] (1.5,2.7) -- (3.5,2.7) node[midway, below, font=\footnotesize] {$\Delta x$};
\draw[|-|, thick, sourceblue] (3.5,2.7) -- (5.5,2.7) node[midway, below, font=\footnotesize] {$\Delta x$};
\draw[|-|, thick, sourceblue] (5.5,2.7) -- (7.5,2.7) node[midway, below, font=\footnotesize] {$\Delta x$};
\draw[|-|, thick, sourceblue] (7.5,2.7) -- (9.5,2.7) node[midway, below, font=\footnotesize] {$\Delta x$};

% Bottom lane: u-ruler (stretched spacing du = g'(x) dx)
\draw[thick, ->] (0.5,0) -- (11.5,0) node[right, font=\small] {$u$};
\node[font=\footnotesize\bfseries, targetgreen!80!black, anchor=north west] at (0.5,-0.75) {Bottom Lane: Stretched $u$-ruler ($\dd u = g'(x)\,\dd x$, where $u=g(x)=x^2$)};

% Transformed ticks on u (e.g. u = 1, 4, 9, 16, 25 mapped to 1.5, 2.5, 4.5, 7.5, 11.0)
\draw[thick, targetgreen!80!black] (1.5,-0.15) -- (1.5,0.15) node[below=4pt, font=\footnotesize] {$u_0{=}1$};
\draw[thick, targetgreen!80!black] (2.5,-0.15) -- (2.5,0.15) node[below=4pt, font=\footnotesize] {$u_1{=}4$};
\draw[thick, targetgreen!80!black] (4.5,-0.15) -- (4.5,0.15) node[below=4pt, font=\footnotesize] {$u_2{=}9$};
\draw[thick, targetgreen!80!black] (7.5,-0.15) -- (7.5,0.15) node[below=4pt, font=\footnotesize] {$u_3{=}16$};
\draw[thick, targetgreen!80!black] (11.0,-0.15) -- (11.0,0.15) node[below=4pt, font=\footnotesize] {$u_4{=}25$};

% Slices on u
\draw[|-|, thick, quotientorange] (1.5,0.4) -- (2.5,0.4) node[midway, above, font=\tiny] {$\Delta u_0 \approx 3$};
\draw[|-|, thick, quotientorange] (2.5,0.4) -- (4.5,0.4) node[midway, above, font=\tiny] {$\Delta u_1 \approx 5$};
\draw[|-|, thick, quotientorange] (4.5,0.4) -- (7.5,0.4) node[midway, above, font=\tiny] {$\Delta u_2 \approx 7$};
\draw[|-|, thick, quotientorange] (7.5,0.4) -- (11.0,0.4) node[midway, above, font=\tiny] {$\Delta u_3 \approx 9$};

% Mapping arrows between rulers through whitespace corridor
\draw[-{Stealth[length=2.5mm]}, dashed, gray!70, thick] (1.5,2.4) to[out=-90, in=90] (1.5,0.8);
\draw[-{Stealth[length=2.5mm]}, dashed, gray!70, thick] (3.5,2.4) to[out=-90, in=90] (2.5,0.8);
\draw[-{Stealth[length=2.5mm]}, dashed, gray!70, thick] (5.5,2.4) to[out=-90, in=90] (4.5,0.8);
\draw[-{Stealth[length=2.5mm]}, dashed, gray!70, thick] (7.5,2.4) to[out=-90, in=90] (7.5,0.8);
\draw[-{Stealth[length=2.5mm]}, dashed, gray!70, thick] (9.5,2.4) to[out=-90, in=90] (11.0,0.8);

\node[font=\footnotesize, fill=white, draw=gray!40, rounded corners=2pt, inner sep=3pt] at (5.5,1.6)
  {Transformation $u = g(x)$: local stretch factor $\dd u / \dd x = g'(x)$};

\end{tikzpicture}
\caption{Substitution as ruler distortion (the 1D Jacobian). \textbf{Top lane:} Slices of equal width $\Delta x$ along the $x$-axis. \textbf{Bottom lane:} Under the mapping $u = g(x)$ (here $u=x^2$), equal steps in $x$ are non-uniformly stretched along the $u$-axis into increments $\Delta u \approx g'(x)\Delta x$. Slices stretch out where $g'(x)$ is large and compress where $g'(x)$ is small. The factor $g'(x)$ converts between rulers so total area is preserved.}
\label{ch22:fig-ruler}
\end{figure}

This is why a substitution can fail or mislead if $g'$ vanishes or changes sign on
the interval: in those places the stretch factor degenerates or flips orientation,
and the naive picture of ``re-measuring widths'' needs care. We will see a clean
example of the sign issue later. For now, hold onto the geometric meaning: $\dd u =
g'(x)\,\dd x$ is the change-of-width law, and substitution is integration performed
in a rescaled coordinate.

\begin{connectionbox}
This change-of-width factor is your first encounter with something that will return
in force later in the book. When you integrate in two or three variables and change
coordinates -- say to polar coordinates -- a similar stretch factor appears, but now
it measures how \emph{areas} (not lengths) get rescaled, and it is called the
\emph{Jacobian}. The single $g'(x)$ here is the one-dimensional Jacobian. Mastering this
one-dimensional stretch factor provides the geometric foundation for the multivariable
change-of-variables formula.
\end{connectionbox}

\section{How to choose $u$}
\label{ch22:choosing-u}

The theorem is easy; the art is picking $u$. What follows is a practical hierarchy
of things to try. None is a guarantee -- integration is a search -- but in that
order they resolve the large majority of substitution problems.

\begin{enumerate}
\item \textbf{Let $u$ be the inside of a composition.} If the integrand contains
something like $\sin(\text{stuff})$, $e^{\text{stuff}}$, $(\text{stuff})^{n}$, or
$\sqrt{\text{stuff}}$, try $u = \text{stuff}$. Then check whether $\dd u =
(\text{stuff})'\,\dd x$ accounts for the rest of the integrand.

\item \textbf{Let $u$ be the thing whose derivative is also present.} Scan for a
chunk of the integrand that is (up to a constant) the derivative of some other
chunk. Name the other chunk $u$. This is the same rule as above, phrased so your eye
hunts for the $g'(x)$ factor.

\item \textbf{Let $u$ be an inconvenient denominator or exponent.} A denominator
$x^2+1$ or an exponent in $e^{3x}$ is often the cleanest choice, especially when its
derivative is lurking in the numerator.

\item \textbf{If a constant is missing, supply it.} You rarely get $\dd u$ exactly;
usually you get a constant multiple of it. Constants are free to move across the
integral sign, so you patch the discrepancy with a constant factor. (Only constants
-- never a variable.)
\end{enumerate}

\begin{trapbox}
You may multiply and divide by \emph{constants} to make $\dd u$ appear, because
$\int c\,h\,\dd x = c\int h\,\dd x$. You may \emph{not} pull a variable across the
integral sign to manufacture $\dd u$. Writing $\int x^2 e^{x}\,\dd x = \tfrac1?\int
\dots$ by ``factoring out an $x$'' is illegal: $\int x\,h(x)\,\dd x \ne x\int
h(x)\,\dd x$. If the leftover factor is not a constant, your chosen $u$ is wrong, or
the integral needs a different technique (such as integration by parts, coming in
the next chapter).
\end{trapbox}

\begin{walkbox}
Before you compute anything, train the eye. Take a walk and, for each of these,
just decide \emph{what you would call $u$} and whether its derivative is present:
$\int x\,e^{x^2}\dd x$; $\int \frac{\ln x}{x}\,\dd x$; $\int \tan x\,\dd x$;
$\int \frac{x}{\sqrt{1-x^2}}\,\dd x$; $\int \cos^3 x\,\dd x$. Three of them are
clean one-step subs, one needs a clever rewrite of $\tan x$, and one needs a trig
identity before the sub appears. Sorting them in your head is the real skill.
\end{walkbox}

\section{Worked examples, from plain to subtle}
\label{ch22:examples}

What follows is a deliberately ordered sequence. Watch how the same single
idea -- recognize $F'(g)\,g'$ -- keeps reappearing in heavier and heavier disguise.

\subsection{The clean cases}

\begin{example}[Inside spotted directly]
\label{ch22:ex-clean1}
Compute $\displaystyle \int 3x^2\,(x^3+1)^5\,\dd x$.

The composition is $(x^3+1)^5$; its inside is $x^3+1$. Let $u = x^3+1$, so
$\dd u = 3x^2\,\dd x$. The integrand contains exactly $3x^2\,\dd x$, which is $\dd
u$. Substituting,
\[
  \int (x^3+1)^5\,\underbrace{3x^2\,\dd x}_{\dd u}
   = \int u^5\,\dd u = \frac{u^6}{6} + C = \frac{(x^3+1)^6}{6} + C.
\]
Differentiate the answer to check: $\ddx{\frac{(x^3+1)^6}{6}} = (x^3+1)^5\cdot 3x^2$.
It matches.
\end{example}

\begin{example}[Patching a missing constant]
\label{ch22:ex-constant}
Compute $\displaystyle \int x\,e^{x^2}\,\dd x$.

Let $u = x^2$, so $\dd u = 2x\,\dd x$. But the integrand has only $x\,\dd x$, not
$2x\,\dd x$. No problem: $x\,\dd x = \tfrac12\,\dd u$. Equivalently, pull the missing
$\tfrac12$ outside:
\[
  \int x\,e^{x^2}\,\dd x = \int e^{u}\cdot \tfrac12\,\dd u
   = \tfrac12\int e^{u}\,\dd u = \tfrac12 e^{u} + C = \tfrac12 e^{x^2} + C.
\]
The half is not a fudge factor. It is the exact bookkeeping the change-of-width law
demands: $\dd x = \tfrac{1}{2x}\,\dd u$, and the stray $x$ in the integrand cancels
the $x$ in $\tfrac{1}{2x}$, leaving the clean $\tfrac12$.
\end{example}

\begin{example}[The logarithmic form]
\label{ch22:ex-log}
Compute $\displaystyle \int \frac{2x}{x^2+1}\,\dd x$.

The denominator $x^2+1$ has derivative $2x$, which is sitting in the numerator. Let
$u = x^2+1$, $\dd u = 2x\,\dd x$:
\[
  \int \frac{2x\,\dd x}{x^2+1} = \int \frac{\dd u}{u} = \ln\abs{u} + C
   = \ln(x^2+1) + C.
\]
(We may drop the absolute value because $x^2+1 > 0$ always.) This is a special case
of the enormously useful pattern $\int \frac{g'(x)}{g(x)}\,\dd x = \ln\abs{g(x)} +
C$: \emph{a fraction whose top is the derivative of its bottom integrates to the log
of the bottom.} Burn that pattern into memory; it is everywhere.
\end{example}

\begin{example}[Two layers, one sub]
\label{ch22:ex-lnx}
Compute $\displaystyle \int \frac{\ln x}{x}\,\dd x$.

Here the ``inside'' is $\ln x$, whose derivative is $\tfrac1x$ -- and there is a
$\tfrac1x$ in the integrand. Let $u = \ln x$, $\dd u = \tfrac1x\,\dd x$:
\[
  \int \ln x\cdot\frac{\dd x}{x} = \int u\,\dd u = \frac{u^2}{2}+C
   = \frac{(\ln x)^2}{2} + C.
\]
\end{example}

\subsection{When the substitution must be rearranged}

Sometimes after choosing $u$ you still have leftover $x$'s that are not part of
$\dd u$. The fix is to solve $u = g(x)$ for $x$ and substitute those leftover
$x$'s too. This is where substitution stops being mere pattern-matching and becomes
a genuine change of variable.

\begin{example}[Solving back for the stray $x$]
\label{ch22:ex-stray}
Compute $\displaystyle \int x\,\sqrt{x+1}\,\dd x$.

The natural choice is $u = x+1$, so $\dd u = \dd x$ and $\sqrt{x+1} = \sqrt u$. But
there is a stray $x$ out front that is not accounted for. Since $u = x+1$, we have
$x = u-1$. Replace it:
\[
  \int x\sqrt{x+1}\,\dd x = \int (u-1)\sqrt u\,\dd u
   = \int \left(u^{3/2} - u^{1/2}\right)\dd u.
\]
Now it is elementary:
\[
  = \frac{2}{5}u^{5/2} - \frac{2}{3}u^{3/2} + C
  = \frac{2}{5}(x+1)^{5/2} - \frac{2}{3}(x+1)^{3/2} + C.
\]
The lesson: a substitution does not have to consume every $x$ through $\dd u$. As
long as you can express \emph{all} remaining $x$'s in terms of $u$, the change of
variable is legitimate and complete.
\end{example}

\begin{example}[A trig identity must come first]
\label{ch22:ex-trigid}
Compute $\displaystyle \int \cos^3 x\,\dd x$.

There is no obvious inside-and-its-derivative here. The move is to spend one factor
of $\cos x$ as the future $\dd u$ and convert the rest using
$\cos^2 x = 1 - \sin^2 x$:
\[
  \int \cos^3 x\,\dd x = \int \cos^2 x\cdot \cos x\,\dd x
   = \int (1-\sin^2 x)\,\cos x\,\dd x.
\]
Now let $u = \sin x$, $\dd u = \cos x\,\dd x$:
\[
  = \int (1-u^2)\,\dd u = u - \frac{u^3}{3} + C
  = \sin x - \frac{\sin^3 x}{3} + C.
\]
The identity created the pattern; substitution finished the job. This trick --
peel off one factor to be $\dd u$, convert the even remainder with a Pythagorean
identity -- is the backbone of trigonometric integrals, which get a chapter of
their own.
\end{example}

\begin{example}[$\tan x$, the classic rewrite]
\label{ch22:ex-tan}
Compute $\displaystyle \int \tan x\,\dd x$.

Rewrite $\tan x = \dfrac{\sin x}{\cos x}$. Now the bottom is $\cos x$ and the top
$\sin x$ is, up to a sign, its derivative. Let $u = \cos x$, $\dd u = -\sin x\,\dd
x$, so $\sin x\,\dd x = -\dd u$:
\[
  \int \frac{\sin x}{\cos x}\,\dd x = \int \frac{-\dd u}{u}
   = -\ln\abs{u} + C = -\ln\abs{\cos x} + C = \ln\abs{\sec x} + C.
\]
This is again the $\int \frac{g'}{g}$ pattern, here applied to trigonometric functions.
\end{example}

\subsection{The genuinely sneaky ones}

\begin{example}[When the obvious sub is not the inside]
\label{ch22:ex-sneaky1}
Compute $\displaystyle \int \frac{1}{x\ln x}\,\dd x$.

You might be tempted by $u = x$ or by the whole denominator. Resist, and notice the
structure: $\tfrac1x$ is the derivative of $\ln x$. So treat $\ln x$ as the
inside. Let $u = \ln x$, $\dd u = \tfrac1x\,\dd x$:
\[
  \int \frac{1}{\ln x}\cdot\frac{\dd x}{x} = \int \frac{\dd u}{u}
   = \ln\abs{u} + C = \ln\abs{\ln x} + C.
\]
The result is a logarithm composed with a logarithm: $\ln\abs{\ln x} + C$.
\end{example}

\begin{example}[Completing the square, then substituting]
\label{ch22:ex-square}
Compute $\displaystyle \int \frac{1}{x^2+4x+5}\,\dd x$.

The denominator does not factor over the reals, but we can complete the square:
$x^2+4x+5 = (x+2)^2 + 1$. Now let $u = x+2$, $\dd u = \dd x$:
\[
  \int \frac{\dd x}{(x+2)^2+1} = \int \frac{\dd u}{u^2+1} = \arctan u + C
   = \arctan(x+2) + C.
\]
Here substitution did something subtler than peel off a chain rule: it
\emph{recentered} the variable so a standard form ($\int \frac{\dd u}{u^2+1} =
\arctan u$) became visible. Shifting and scaling are the humblest and most common
substitutions of all.
\end{example}

\begin{example}[A rationalizing substitution]
\label{ch22:ex-rational}
Compute $\displaystyle \int \frac{1}{1+\sqrt{x}}\,\dd x$.

The square root is the obstruction, so give it a name: let $u = \sqrt x$, which means
$x = u^2$ and $\dd x = 2u\,\dd u$. The integral becomes
\[
  \int \frac{2u\,\dd u}{1+u}.
\]
Now rewrite the numerator to split off a whole copy of the denominator:
$\dfrac{2u}{1+u} = \dfrac{2(u+1) - 2}{1+u} = 2 - \dfrac{2}{1+u}$. Hence
\[
  \int\left(2 - \frac{2}{1+u}\right)\dd u = 2u - 2\ln\abs{1+u} + C
   = 2\sqrt x - 2\ln\!\big(1+\sqrt x\big) + C.
\]
Notice this substitution ran ``the other way'': instead of setting $u$ equal to a
chunk of the integrand and computing $\dd u$, we set $x$ equal to a function of $u$
and computed $\dd x$. Both directions are legal; Theorem \ref{ch22:thm-indef}
covers them as long as the change of variable is invertible on the interval in
question. (Here $u=\sqrt x \ge 0$, which is fine for $x \ge 0$.)
\end{example}

\begin{aestheticbox}
Look back at how many different-looking integrals collapsed onto the same few
shapes: $\int u^n\,\dd u$, $\int e^u\,\dd u$, $\int \frac{\dd u}{u}$, and $\int
\frac{\dd u}{u^2+1}$. Substitution is, in a sense, a \emph{universal translator}: it
keeps trying to phrase the integral in one of the handful of languages you already
speak fluently. The art of integration is largely the art of recognizing which
translation will land you in familiar territory.
\end{aestheticbox}

\section{Definite integrals: carry the limits}
\label{ch22:definite}

For a definite integral you have two equally valid styles, and you should be fluent
in both.

\textbf{Style A -- change the limits.} Use Theorem \ref{ch22:thm-def} directly:
when $u = g(x)$, replace the limits $x=a, x=b$ by $u=g(a), u=g(b)$ and never return
to $x$ at all.

\textbf{Style B -- back-substitute.} Do the indefinite integral, return to $x$,
then plug in the original limits.

\begin{example}[The same integral both ways]
\label{ch22:ex-bothways}
Compute $\displaystyle \int_0^2 x\,e^{x^2}\,\dd x$.

Let $u = x^2$, $\dd u = 2x\,\dd x$, so $x\,\dd x = \tfrac12\,\dd u$.

\emph{Style A (change the limits).} When $x=0$, $u=0$; when $x=2$, $u=4$. So
\[
  \int_0^2 x\,e^{x^2}\,\dd x = \frac12\int_{0}^{4} e^u\,\dd u
   = \frac12\eval{e^u}{0}^{4} = \frac{e^4-1}{2}.
\]
\emph{Style B (back-substitute).} The antiderivative is $\tfrac12 e^{x^2}$, so
\[
  \int_0^2 x\,e^{x^2}\,\dd x = \frac12\eval{e^{x^2}}{0}^{2}
   = \frac12\big(e^4 - e^0\big) = \frac{e^4-1}{2}.
\]
Same answer, as it must be. Style A is usually faster and cleaner once you trust it,
because you never untangle $u$ back into $x$.
\end{example}

\begin{recoverybox}
\textbf{Pause and think.} Suppose you apply the substitution $u = 3x+1$ to the definite integral $\int_1^5 f(3x+1)\,\dd x$. What are the new limits of integration in $u$, before you know anything about $f$?

\emph{Answer:} When $x = 1$, $u = 3(1)+1 = 4$. When $x = 5$, $u = 3(5)+1 = 16$. The integral becomes $\frac{1}{3}\int_4^{16} f(u)\,\dd u$.
\end{recoverybox}

\begin{trapbox}
The signature blunder: changing the variable to $u$ but leaving the old $x$-limits
on the integral sign. If you write
\[
  \int_0^2 x\,e^{x^2}\,\dd x \;\overset{?}{=}\; \frac12\int_0^2 e^u\,\dd u,
\]
you have committed it -- the ``$0$'' and ``$2$'' there are $x$-values, but the
integral is now in $u$. Either change them to the $u$-values ($0$ and $4$), or drop
the limits entirely until you have returned to $x$. Mixing $u$-integrand with
$x$-limits gives a wrong number every time. A safe habit: the instant you write a
$u$ on the integrand, write the $u$-limits on the sign, in the same stroke of the
pen.
\end{trapbox}

\begin{figure}[H]
\centering
\begin{tikzpicture}
\begin{axis}[
  name=plotx,
  width=6.6cm, height=5.5cm,
  title={$x$-plane: $\int_0^2 x e^{x^2}\,\dd x$},
  title style={font=\small},
  xlabel={$x$}, ylabel={$y$},
  domain=0:2, samples=100,
  axis lines=left,
  xmin=0, xmax=2.15,
  ymin=0, ymax=118,
  enlargelimits=false,
]
\addplot[curveblue, thick, name path=fx] {x*exp(x^2)};
\path[name path=axisx] (axis cs:0,0) -- (axis cs:2,0);
\addplot[areafill!60] fill between[of=fx and axisx, soft clip={domain=0:2}];
\node[font=\footnotesize, fill=white, inner sep=2pt, rounded corners=2pt] at (axis cs:1.1,40) {Area $\approx 26.8$};
\end{axis}

\begin{axis}[
  at={($(plotx.east)+(1.1cm,0)$)},
  anchor=west,
  width=6.6cm, height=5.5cm,
  title={$u$-plane: $\frac{1}{2}\int_0^4 e^u\,\dd u$},
  title style={font=\small},
  xlabel={$u$}, ylabel={$y$},
  domain=0:4, samples=100,
  axis lines=left,
  xmin=0, xmax=4.3,
  ymin=0, ymax=60,
  enlargelimits=false,
]
\addplot[targetgreen!80!black, thick, name path=fu] {0.5*exp(x)};
\path[name path=axisu] (axis cs:0,0) -- (axis cs:4,0);
\addplot[targetgreen!30] fill between[of=fu and axisu, soft clip={domain=0:4}];
\node[font=\footnotesize, fill=white, inner sep=2pt, rounded corners=2pt] at (axis cs:2.2,20) {Area $\approx 26.8$};
\end{axis}
\end{tikzpicture}
\caption{Preservation of area under substitution. \textbf{Left:} In $x$-space, the integral $\int_0^2 x e^{x^2}\,\dd x$ forms a tall, narrow region of width $2$ and peak height $\approx 109.2$. \textbf{Right:} The substitution $u = x^2$ stretches the base onto $[0,4]$ while the factor $\dd x = \frac{1}{2x}\dd u$ scales down the height to $\frac{1}{2}e^u$ (peak $\approx 54.6$). The transformed region has the exact same area, $\frac{e^4-1}{2} \approx 26.8$. Changing the limits from $x\in[0,2]$ to $u\in[0,4]$ is what ensures the re-measured area stays equal.}
\label{ch22:fig-area}
\end{figure}

\begin{labbox}
Verify the limit-changing rule on a case where you know the
answer directly. Take $\int_1^3 (2x)(x^2)\,\dd x$, which is just $\int_1^3 2x^3\,\dd x =
\tfrac12 x^4 \big|_1^3 = \tfrac12(81-1) = 40$. Now redo it by substituting $u =
x^2$, $\dd u = 2x\,\dd x$, with limits $u=1$ to $u=9$: you get $\int_1^9
u\,\dd u = \tfrac12 u^2\big|_1^9 = \tfrac12(81-1) = 40$. The limits change from
$\{1,3\}$ to $\{1,9\}$, while the evaluated area remains $40$, exactly as guaranteed by the definite-form theorem.
\end{labbox}

\subsection{A symmetry payoff}

Changing limits is not just bookkeeping; sometimes it reveals structure you could
not otherwise see.

\begin{example}[Substitution exposes a symmetry]
\label{ch22:ex-symmetry}
Show that $\displaystyle \int_0^{\pi/2} \frac{\sin x}{\sin x + \cos x}\,\dd x =
\frac{\pi}{4}$.

There is no friendly antiderivative here. Instead, substitute $u = \tfrac{\pi}{2} -
x$, the reflection that swaps the endpoints. Then $\dd u = -\dd x$; the limits flip
($x=0\!\to\!u=\tfrac\pi2$ and $x=\tfrac\pi2\!\to\!u=0$); and $\sin x = \cos u$,
$\cos x = \sin u$. Call the original integral $I$:
\[
  I = \int_0^{\pi/2}\frac{\sin x}{\sin x+\cos x}\,\dd x
    = \int_{\pi/2}^0 \frac{\sin(\pi/2-u)}{\sin(\pi/2-u)+\cos(\pi/2-u)}\,(-\dd u)
    = \int_0^{\pi/2}\frac{\cos u}{\cos u + \sin u}\,\dd u.
\]
The negative sign on $-\dd u$ reverses the reversed limits $\int_{\pi/2}^0$ back to
their standard order $\int_0^{\pi/2}$. The resulting integral has the exact same value
as the original, but with $\sin$ and $\cos$ traded. Now add the two expressions for $I$:
\[
  2I = \int_0^{\pi/2}\frac{\sin x + \cos x}{\sin x + \cos x}\,\dd x
     = \int_0^{\pi/2} 1\,\dd x = \frac{\pi}{2},
\]
so $I = \tfrac{\pi}{4}$. A substitution did what no antiderivative could: it turned
an unsolvable-looking integral into a triviality by exploiting symmetry. Figure
\ref{ch22:fig-symmetry} shows why: the integrand and its reflection add to the
constant $1$, so their two equal areas together fill a rectangle of area
$\tfrac\pi2$.
\end{example}

\begin{figure}[H]
\centering
\begin{tikzpicture}
\begin{axis}[
  width=12cm, height=6cm,
  xlabel={$x$}, ylabel={$y$},
  domain=0:1.5708, samples=120,
  axis lines=left,
  ymin=0, ymax=1.18,
  xmin=0, xmax=1.5708,
  xtick={0,0.7854,1.5708},
  xticklabels={$0$,$\pi/4$,$\pi/2$},
  enlargelimits=false,
  legend style={at={(0.5,1.03)}, anchor=south, legend columns=3, font=\footnotesize, /tikz/every even column/.append style={column sep=12pt}},
  legend cell align=left,
]
\path[name path=axis] (axis cs:0,0) -- (axis cs:1.5708,0);
\addplot[curveblue, thick, name path=f] {sin(deg(x))/(sin(deg(x))+cos(deg(x)))};
\addlegendentry{$\tfrac{\sin x}{\sin x+\cos x}$}
\addplot[areafill!45, forget plot] fill between[of=f and axis, soft clip={domain=0:1.5708}];
\addplot[tangentred, thick, dashed] {cos(deg(x))/(sin(deg(x))+cos(deg(x)))};
\addlegendentry{$\tfrac{\cos x}{\sin x+\cos x}$}
\addplot[black, dotted, thick] {1};
\addlegendentry{their sum $=1$}
\node[font=\footnotesize, sourceblue, fill=white, inner sep=2pt, rounded corners=1pt] at (axis cs:1.15,0.3) {Area $I = \frac{\pi}{4}$};
\node[font=\footnotesize, tangentred, fill=white, inner sep=2pt, rounded corners=1pt] at (axis cs:0.42,0.7) {Area $I = \frac{\pi}{4}$};
\end{axis}
\end{tikzpicture}
\caption{The original integrand (solid) and its reflection $x\mapsto\tfrac{\pi}{2}-x$
(dashed) are mirror images across $x=\pi/4$, and at every point they sum to $1$
(dotted). The two shaded regions each have area $I = \pi/4$ and together fill the
rectangle $[0, \pi/2]\times[0,1]$ of area $\pi/2$. Substitution proves their equality without evaluating an antiderivative.}
\label{ch22:fig-symmetry}
\end{figure}

\begin{connectionbox}
This reflection identity generalizes to the property that $\int_a^b f(x)\,\dd x =
\int_a^b f(a+b-x)\,\dd x$ for any integrable $f$ -- a direct consequence of the
substitution $u = a+b-x$. It is widely used in problem solving and in
probability, where it underlies calculations with symmetric distributions.
\end{connectionbox}

\section{What the symbol-shuffling really means}
\label{ch22:meaning}

Time to face head-on the thing that should still nag a careful reader. When we
write $\dd u = g'(x)\,\dd x$ and then ``cancel'' $\dd x$'s, are we doing real
mathematics or performing a ritual that happens to work?

\begin{happeningbox}
It is real mathematics, but the symbols are doing something subtler than the ritual
suggests. Strictly, $\dd u = g'(x)\,\dd x$ is \emph{defined} to be shorthand for the
statement of Theorem \ref{ch22:thm-def}. The $\dd x$ standing alone is not a number,
so ``cancelling'' it is not division. What justifies the manipulation is that the
shorthand was \emph{designed} to mirror a true theorem: every legal token-move you
make with the $\dd$'s corresponds to a legal step in the proof.

Two things keep you safe. First, the substitution $u=g(x)$ should be
\emph{differentiable}, so that $g'(x)$ exists and the change-of-width law makes
sense. Second, for definite integrals you want $g'$ to keep a constant sign (or you
must track orientation carefully), and for the reverse-direction substitutions
$x=h(u)$ of Example \ref{ch22:ex-rational} you want $h$ to be \emph{invertible} on
the interval -- meaning one-to-one, so every output comes from exactly one input and
the trip back to $x$ is unambiguous. When those mild conditions hold, the
symbol-shuffling is guaranteed correct, because it is just a compressed transcription
of the theorem we proved in three lines.
\end{happeningbox}

\begin{trapbox}
The conditions are not decorative. Consider $\int_{-1}^{1} x^2\,\dd x = \tfrac23$,
and try the substitution $u = x^2$. It is tempting to write $\dd u = 2x\,\dd x$, so
$x\,\dd x = \tfrac12\dd u$ -- but the integrand is $x^2\,\dd x$, not $x\,\dd x$, and
solving $x = \pm\sqrt u$ forces a sign \emph{choice} that is positive on $[0,1]$ and
negative on $[-1,0]$. The map $x\mapsto x^2$ is not invertible across $0$ (it folds
the interval in half), so a single substitution over the whole interval is illegal.
The fix is to split at $x=0$ and substitute on each piece separately. Moral:
when $g'$ changes sign -- here at $x=0$ -- the change of variable folds the domain
and quietly double-counts or cancels. Split first, then substitute.
\end{trapbox}

\begin{figure}[H]
\centering
\begin{tikzpicture}
\begin{axis}[
  width=10.5cm, height=5.8cm,
  xlabel={$x$}, ylabel={$u = g(x) = x^2$},
  domain=-1:1, samples=100,
  axis lines=middle,
  xmin=-1.35, xmax=1.35,
  ymin=-0.22, ymax=1.18,
  xtick={-1, 1},
  ytick={0, 1},
  legend style={at={(0.5,1.03)}, anchor=south, legend columns=2, font=\footnotesize, /tikz/every even column/.append style={column sep=14pt}},
  legend cell align=left,
]
\addplot[quotientorange, very thick, domain=-1:0] {x^2};
\addlegendentry{Left branch: $g'(x) < 0$}
\addplot[curveblue, very thick, domain=0:1] {x^2};
\addlegendentry{Right branch: $g'(x) > 0$}

\draw[dashed, gray!80, thick] (axis cs:-0.7071,0) -- (axis cs:-0.7071,0.5) -- (axis cs:0.7071,0.5) -- (axis cs:0.7071,0);
\node[circle, fill=quotientorange, inner sep=1.8pt] at (axis cs:-0.7071,0.5) {};
\node[circle, fill=curveblue, inner sep=1.8pt] at (axis cs:0.7071,0.5) {};
\node[above, font=\footnotesize, fill=white, inner sep=2pt, rounded corners=1pt] at (axis cs:0,0.5) {single $u = 0.5$};
\node[below, font=\footnotesize, quotientorange, fill=white, inner sep=1.5pt] at (axis cs:-0.7071,-0.03) {$-\frac{1}{\sqrt{2}}$};
\node[below, font=\footnotesize, curveblue, fill=white, inner sep=1.5pt] at (axis cs:0.7071,-0.03) {$+\frac{1}{\sqrt{2}}$};
\node[right, font=\footnotesize, fill=white, inner sep=2pt, rounded corners=1pt] at (axis cs:0.06,0.12) {fold point: $g'(0)=0$};
\end{axis}
\end{tikzpicture}
\caption{Why sign changes in $g'$ break a single substitution. The mapping $u = x^2$
folds the interval $[-1,1]$ in half at the stationary point $x=0$. Both the left branch
(orange, where $g'<0$) and the right branch (blue, where $g'>0$) map onto the same range
$u\in[0,1]$, so a single $u$-value corresponds to two distinct $x$-values. To substitute safely,
split the integral at $x=0$ so $g$ is one-to-one on each piece.}
\label{ch22:fig-folding}
\end{figure}

\begin{intuitionbox}
You can always check an antiderivative directly by differentiating it: if
$\ddx{(\text{your answer})}$ recovers the original integrand, your answer is right.
\end{intuitionbox}

\begin{applicationbox}
Substitution is not a classroom curiosity; it is how real models get evaluated.
\emph{Physics:} the work done by a varying force, $\int F\,\dd x$, becomes a clean
energy integral after substituting the position-dependent variable. \emph{Probability
and ML:} normalizing a Gaussian relies on the substitution $u = x/\sigma$ to reduce
$\int e^{-x^2/2\sigma^2}\dd x$ to the standard bell curve; the same change of
variable is the ``reparameterization trick'' that lets modern generative models
backpropagate through random sampling. \emph{Biology and chemistry:} separable
differential equations for population growth and reaction rates are solved by
substituting to isolate each variable, then integrating. The little factor $g'(x)$
you have been tracking is, in every one of these, a unit conversion between the
natural variable of the problem and the variable that makes the integral solvable.
\end{applicationbox}

\section{A compact procedure}
\label{ch22:procedure}

Here is the chapter distilled into a checklist you can run on autopilot.

\begin{enumerate}
\item \textbf{Spot the inside.} Find a composition and let $u$ be its inside (or:
find a chunk whose derivative also appears, and call that chunk $u$).
\item \textbf{Differentiate:} compute $\dd u = g'(x)\,\dd x$.
\item \textbf{Rewrite the whole integrand} in terms of $u$ and $\dd u$. Patch
missing \emph{constants} freely; if a stray \emph{variable} remains, solve $u=g(x)$
for $x$ and substitute it too.
\item \textbf{For a definite integral, change the limits} to $u=g(a)$ and
$u=g(b)$ at the same moment you change the integrand.
\item \textbf{Integrate in $u$}, landing in one of your known forms.
\item \textbf{Translate back to $x$} (indefinite) or just evaluate at the
$u$-limits (definite).
\item \textbf{Differentiate your answer to check.} Always.
\end{enumerate}

\begin{recoverybox}
If a substitution leaves you with a leftover variable you cannot express in $u$, do
not force it. This indicates either that a different choice of $u$ is needed, or
that the integral requires another method. The next chapter introduces \emph{integration by
parts}, the technique for integrands that are products with no inner-derivative
factor -- the natural complement to substitution.
\end{recoverybox}

\section{Limitations and the road ahead}

\begin{connectionbox}
Substitution is powerful, but it has a clear limitation: it requires a matching
derivative factor $g'(x)$ inside the integrand (up to an adjustable constant).

When that factor is present, as in $\int x e^{x^2}\,\dd x$, the method works directly:
setting $u = x^2$ gives $\dd u = 2x\,\dd x$, turning the integral into
$\frac{1}{2}\int e^u\,\dd u = \frac{1}{2}e^{x^2} + C$.

When that factor is missing, substitution fails. Compare $\int x e^{x^2}\,\dd x$ with
$\int x e^x\,\dd x$. In $\int x e^x\,\dd x$, setting $u = x$ changes nothing, while
setting $u = e^x$ yields $\int \ln(u)\,\dd u$. Neither factor is the derivative of an
internal function. Similar difficulties arise with integrals such as $\int \ln x\,\dd x$,
$\int x^2\cos x\,\dd x$, and $\int \arctan x\,\dd x$.

To handle products of two unrelated functions, we need a method that inverts the
Product Rule rather than the Chain Rule. That method is \emph{integration by parts},
developed in Chapter~23:
\[
  \int u\,\dd v = uv - \int v\,\dd u.
\]
\end{connectionbox}

\section{Exercises}
\label{ch22:exercises}

\begin{exercise}[Mental checks: spot the pair in five seconds]
For each integral, identify the inside function $u$ and the derivative factor $g'(x)\,\dd x$ without pencil and paper, and state the antiderivative:
\begin{enumerate}
\item $\displaystyle \int 2x\left(x^2+9\right)^4\,\dd x$
\item $\displaystyle \int \frac{\cos x}{\sin x}\,\dd x$
\item $\displaystyle \int \frac{(\ln x)^3}{x}\,\dd x$
\item $\displaystyle \int e^{\tan x}\sec^2 x\,\dd x$
\end{enumerate}
\end{exercise}

\begin{exercise}[Warm-ups]
Evaluate each indefinite integral by a single substitution, and check by
differentiating.
\begin{enumerate}
\item $\displaystyle \int (2x+3)^7\,\dd x$
\item $\displaystyle \int x^2\sqrt{x^3+1}\,\dd x$
\item $\displaystyle \int \frac{e^{1/x}}{x^2}\,\dd x$
\item $\displaystyle \int \sin x\,\cos^4 x\,\dd x$
\end{enumerate}
\end{exercise}

\begin{exercise}[The $g'/g$ pattern]
Using $\int \frac{g'}{g}\,\dd x = \ln\abs{g} + C$, evaluate
$\displaystyle \int \frac{3x^2 - 2}{x^3 - 2x + 7}\,\dd x$ and
$\displaystyle \int \cot x\,\dd x$. For the second, first rewrite $\cot x$ as a
quotient.
\end{exercise}

\begin{exercise}[Stray variable]
Evaluate $\displaystyle \int x(2x-1)^{10}\,\dd x$ by letting $u = 2x-1$ and solving
for the leftover $x$.
\end{exercise}

\begin{exercise}[Definite, change the limits]
Evaluate $\displaystyle \int_0^{1} \frac{x}{x^2+1}\,\dd x$ two ways: by changing the
limits, and by back-substituting. Confirm the answers agree, and state the new
limits explicitly.
\end{exercise}

\begin{exercise}[Complete the square]
Evaluate $\displaystyle \int \frac{1}{x^2 - 6x + 13}\,\dd x$ by completing the square
and substituting to reach an $\arctan$.
\end{exercise}

\begin{exercise}[Rationalize]
Evaluate $\displaystyle \int \frac{1}{1+\sqrt[3]{x}}\,\dd x$ by setting $u =
\sqrt[3]{x}$. (You will need polynomial division.)
\end{exercise}

\begin{exercise}[A trig identity first]
Evaluate $\displaystyle \int \sin^2 x\,\cos^3 x\,\dd x$. (Peel off one $\cos x$ for
$\dd u$ and convert the rest with $\cos^2 x = 1-\sin^2 x$.)
\end{exercise}

\begin{exercise}[Symmetry]
Use the reflection substitution $u = \tfrac{\pi}{2}-x$ to evaluate the definite integral
$\displaystyle \int_0^{\pi/2} \frac{\cos^2 x}{\sin^2 x + \cos^2 x}\,\dd x$ without
finding an antiderivative. (It should remind you of Example \ref{ch22:ex-symmetry}.)
\end{exercise}

\begin{exercise}[Conceptual: why the constant is allowed]
Explain in a sentence or two why you are permitted to pull a constant factor out of
an integral to manufacture $\dd u$, but \emph{not} a variable factor. Illustrate
with a wrong ``computation'' of $\int x^2 e^{x}\,\dd x$ that mistakenly factors out
an $x$, and point to the exact illegal step.
\end{exercise}

\begin{exercise}[Conceptual: the folded interval]
Consider $\displaystyle \int_{-1}^{2} x^2\,\dd x$ and the substitution $u = x^2$.
Explain why a naive single substitution over $[-1,2]$ is invalid, identify where
$g'(x)$ changes sign, and describe how you would split the interval to substitute
correctly. (You do not need to finish the computation.)
\end{exercise}

\begin{exercise}[Proof]
Prove the reflection identity $\displaystyle \int_a^b f(x)\,\dd x = \int_a^b
f(a+b-x)\,\dd x$ for continuous $f$, by carrying out the substitution $u = a+b-x$
carefully, including what happens to the limits and the sign of $\dd u$. Then use it
to evaluate $\int_0^1 \frac{1}{1 + e^{\,x - 1/2}}\,\dd x$ -- or, more simply, re-derive Example \ref{ch22:ex-symmetry} as a corollary.
\end{exercise}

\begin{exercise}[Proof: substitution from the Riemann sum]
For a strictly increasing, continuously differentiable $g$ on $[a,b]$, sketch why
the change-of-variable formula holds at the level of Riemann sums: explain how a
partition of the $x$-axis maps to a partition of the $u$-axis, why the width of the
$k$-th subinterval transforms by approximately the factor $g'(x_k)$, and how letting
the mesh go to zero recovers $\int_{g(a)}^{g(b)} f(u)\,\dd u$. This is the
``change-of-width law'' of Section \ref{ch22:what-dx-does} made into a proof sketch.
\end{exercise}
